% Options for packages loaded elsewhere
\PassOptionsToPackage{unicode}{hyperref}
\PassOptionsToPackage{hyphens}{url}
\PassOptionsToPackage{dvipsnames,svgnames,x11names}{xcolor}
\documentclass[
  10pt,
  letterpaper,
]{article}
\usepackage{xcolor}
\usepackage[margin=1in]{geometry}
\usepackage{amsmath,amssymb}
\setcounter{secnumdepth}{-\maxdimen} % remove section numbering
\usepackage{iftex}
\ifPDFTeX
  \usepackage[T1]{fontenc}
  \usepackage[utf8]{inputenc}
  \usepackage{textcomp} % provide euro and other symbols
\else % if luatex or xetex
  \usepackage{unicode-math} % this also loads fontspec
  \defaultfontfeatures{Scale=MatchLowercase}
  \defaultfontfeatures[\rmfamily]{Ligatures=TeX,Scale=1}
\fi
\usepackage{lmodern}
\ifPDFTeX\else
  % xetex/luatex font selection
\fi
% Use upquote if available, for straight quotes in verbatim environments
\IfFileExists{upquote.sty}{\usepackage{upquote}}{}
\IfFileExists{microtype.sty}{% use microtype if available
  \usepackage[]{microtype}
  \UseMicrotypeSet[protrusion]{basicmath} % disable protrusion for tt fonts
}{}
\makeatletter
\@ifundefined{KOMAClassName}{% if non-KOMA class
  \IfFileExists{parskip.sty}{%
    \usepackage{parskip}
  }{% else
    \setlength{\parindent}{0pt}
    \setlength{\parskip}{6pt plus 2pt minus 1pt}}
}{% if KOMA class
  \KOMAoptions{parskip=half}}
\makeatother
\setlength{\emergencystretch}{3em} % prevent overfull lines
\providecommand{\tightlist}{%
  \setlength{\itemsep}{0pt}\setlength{\parskip}{0pt}}
\usepackage{bookmark}
\IfFileExists{xurl.sty}{\usepackage{xurl}}{} % add URL line breaks if available
\urlstyle{same}
\hypersetup{
  colorlinks=true,
  linkcolor={Maroon},
  filecolor={Maroon},
  citecolor={Blue},
  urlcolor={blue},
  pdfcreator={LaTeX via pandoc}}

\author{}
\date{}

\begin{document}

\section{Response to the reviewer}\label{response-to-the-reviewer}

\textbf{Manuscript:} \emph{Leakage-Safe Multi-Seed Evaluation of
Published BU-Net Components for Resource-Constrained 2D Glioma
Segmentation}

We thank the reviewer for identifying problems that could not be
corrected by editorial changes alone. We rebuilt the study around a
canonical cohort, patient-level partitions, matched training, guarded
test access, reproducible artifacts, and appropriately limited claims.
The old cross-architecture ranking and clinical or state-of-the-art
language were removed.

\subsection{1. Attribution and novelty}\label{attribution-and-novelty}

\textbf{Concern:} RES and WC appeared to be presented as new
contributions although they were introduced in BU-Net.

\textbf{Response:} Corrected throughout. RES, WC, and their combination
are explicitly attributed to Rehman et al.~(2020). The title, abstract,
architecture figure, Methods, Discussion, and Conclusion describe this
as an evaluation and reimplementation study. No architectural novelty is
claimed. A dedicated implementation note records deliberate deviations
from the source paper.

\subsection{2. BraTS 2019/2020 overlap and
leakage}\label{brats-20192020-overlap-and-leakage}

\textbf{Concern:} Pooling BraTS editions could duplicate patients and
contaminate splits.

\textbf{Response:} We did not pool the editions. BraTS 2020 training is
the only canonical labeled cohort. BraTS 2019 is used only for
content-based identity auditing. All 335 BraTS 2019 cases map to BraTS
2020; 34 cases are new in 2020. One mapped segmentation contains a
two-voxel annotation revision. The final 369 patients were split at
patient level into 258/37/74 training/validation/internal-test cases,
with zero identifier or content-signature overlap.

\subsection{3. Slice-level leakage}\label{slice-level-leakage}

\textbf{Concern:} A slice-level split would allow images from one
patient in several subsets.

\textbf{Response:} All partitions and statistical analyses now use the
patient as the unit. Slices are sampled only within an already assigned
training patient. The split manifests and hashes are frozen and tested
automatically.

\subsection{4. Unfair or inconsistent
protocols}\label{unfair-or-inconsistent-protocols}

\textbf{Concern:} Architectures appeared to receive different losses,
schedules, or training opportunities.

\textbf{Response:} Every final candidate now uses the same inputs,
normalization, augmentation, BCE+FTL objective, AdamW optimizer,
learning rate, weight decay, batch size, scheduler, device, and
2,000-step/0.5-hour cap. The protocol is intentionally bounded and this
limitation is stated. Development screens, confirmation runs, and test
evaluation have distinct roles and machine-readable registries.

\subsection{5. Unsupported transformer or Swin
conclusions}\label{unsupported-transformer-or-swin-conclusions}

\textbf{Concern:} Undertrained transformer comparisons cannot support
architectural conclusions.

\textbf{Response:} Removed. No Swin, transformer, attention, or
literature score appears as an experimental comparator. nnU-Net,
TransBTS, and nnFormer are cited only as important untested context. The
manuscript explicitly states that no 3D, transformer, self-configuring,
or external-validation comparison was performed.

\subsection{6. Missing ablation}\label{missing-ablation}

\textbf{Concern:} The contribution of individual components was unclear.

\textbf{Response:} The rebuilt design screens standard U-Net, U-Net+RES,
U-Net+WC, BU-Net (RES+WC), residual-block U-Net, and residual-block
U-Net+WC under one protocol. The final frozen comparison includes
standard U-Net, U-Net+RES, and BU-Net. On the internal test, mean
regional Dice is 0.736, 0.756, and 0.752, respectively. The paired
U-Net+RES versus standard U-Net difference is +0.020 (Holm p=0.004); the
BU-Net versus standard U-Net difference is +0.017 (Holm p=0.014). We
interpret this as evidence for RES under the present protocol, not
architectural novelty.

\subsection{7. Loss definition and
selection}\label{loss-definition-and-selection}

\textbf{Concern:} The loss formulation, class handling, and selection
process were insufficiently specified.

\textbf{Response:} The Methods now provide the complete BCE+FTL formula
and its parameters: equal term weights, alpha 0.3, beta 0.7, gamma 0.75,
smoothing 1e-5, foreground-only FTL, and no class weights. Seven loss
arms were screened once on the validation subset. Their full results are
reported as development-only evidence.

\subsection{8. Inconsistent metrics and
aggregation}\label{inconsistent-metrics-and-aggregation}

\textbf{Concern:} Metric definitions, empty-mask rules, and aggregation
could change model rankings.

\textbf{Response:} The evaluator now fixes WT/TC/ET label mapping, Dice,
HD95, surface Dice at 1 mm, lesion metrics, 26-connectivity,
maximum-total-IoU lesion matching, and explicit empty-mask behavior.
Seeds are averaged within patient before inference; patients are never
multiplied into slice- or seed-level pseudo-samples. Infinite and
undefined secondary values are retained and counted.

\subsection{9. Statistics and
uncertainty}\label{statistics-and-uncertainty}

\textbf{Concern:} Point estimates without patient-level uncertainty or
multiplicity control were inadequate.

\textbf{Response:} The primary outcome and three contrasts were frozen
before test access. We use 10,000 patient-level bootstrap resamples,
100,000 two-sided paired sign-flip permutations, paired effect size dz,
and Holm correction. Regional and secondary endpoints are identified as
estimation-only.

\subsection{10. Resource and efficiency
claims}\label{resource-and-efficiency-claims}

\textbf{Concern:} ``Lightweight,'' ``efficient,'' or clinical deployment
language lacked measurement.

\textbf{Response:} We removed those unqualified claims. The paper
reports parameters, checkpoint size, MACs, FLOPs, p50/p95 volume
latency, throughput, allocated/reserved memory, and development
accelerator-hours from the implemented candidates on one host. BU-Net is
explicitly shown to require more resources than standard U-Net and
U-Net+RES. No clinical deployment claim remains.

\subsection{11. Qualitative analysis and failure
cases}\label{qualitative-analysis-and-failure-cases}

\textbf{Concern:} The paper lacked image-level evidence and failure
analysis.

\textbf{Response:} Artifact-derived figures now show modalities, ground
truth, all frozen predictions, false-positive/false-negative overlays,
and preselected success, hard, and failure cases. The failure case is
discussed alongside modest ET lesion-recall and lesion-wise Dice
estimates.

\subsection{12. Reproducibility and
code}\label{reproducibility-and-code}

\textbf{Concern:} The reported results could not be reproduced.

\textbf{Response:} The repository now contains configurations, model and
loss implementations, evaluator tests, patient manifests and hashes, run
registries, statistical scripts, derived tables and figures, environment
lock, claim ledger, and reproduction instructions. A clean-clone audit
verified 230 tracked-artifact hashes, static checks, the test suite, two
byte-identical manuscript-input generations, and a clean worktree.

\subsection{13. Requested 3D, nnU-Net, transformer, and external
validation}\label{requested-3d-nnu-net-transformer-and-external-validation}

\textbf{Concern:} Strong conclusions would require modern
3D/self-configuring baselines and external validation.

\textbf{Response:} We agree with the evidentiary requirement and
narrowed the paper instead of fabricating or importing incomparable
results. The revised manuscript does not claim superiority over these
systems, state of the art, external generalization, or clinical
applicability. These experiments remain future work and are prominent
limitations.

\subsection{14. Manuscript-wide claim
correction}\label{manuscript-wide-claim-correction}

\textbf{Concern:} The original framing was broader than the evidence.

\textbf{Response:} The manuscript was rewritten. The conclusion is
limited to a leakage-safe internal comparison: RES improved the primary
overlap endpoint over standard U-Net under this bounded 2D protocol;
adding WC in the full BU-Net did not improve that endpoint over RES
alone and increased resource demand. Boundary and lesion-level findings
are reported as mixed.

\end{document}
